\chapter{Radiation: RRTMG longwave and Dudhia shortwave}\label{ch:rad}

Radiation heats and cools the atmosphere and drives the surface energy balance. The case uses:
\begin{itemize}
  \item the longwave part of the Rapid Radiative Transfer Model for GCMs, RRTMG \citep{mlawer1997,iacono2008}:
    \code{ra_lw_physics = 4}, \code{phys/module_ra_rrtmg_lw.F};
  \item the shortwave scheme of \citet{dudhia1989}: \code{ra_sw_physics = 1}, \code{phys/module_ra_sw.F}.
\end{itemize}
Both are called from \code{phys/module_radiation_driver.F} every \code{radt} minutes. Between calls, the heating
rates are kept and applied every step.

\section{What the driver does}

On every step, \src{phys/module_radiation_driver.F}{radiation_driver} accumulates the surface fluxes for output and
computes the total cloud cover of each column ($1 - \prod_k(1 - c_k)$, in $k$ order). On radiation steps it also:
\begin{enumerate}
  \item computes the solar declination and the cosine of the zenith angle at every point (\code{radconst},
    \code{calc_coszen}), and the eclipse obscuration (zero for this case);
  \item diagnoses the cloud fraction (\code{cal_cldfra1}) from the cloud water and ice;
  \item interpolates the ozone climatology in time (\code{ozn_time_int}) and in pressure (\code{ozn_p_int}), on the
    outer domain;
  \item calls RRTMG longwave, then Dudhia shortwave;
  \item splits the downward shortwave into direct and diffuse parts, and stores the heating rates.
\end{enumerate}

\section{RRTMG longwave: the correlated-k method}

A line-by-line calculation of thermal radiation would integrate over hundreds of thousands of absorption lines. The
correlated-$k$ method replaces the integration over wavenumber within a band by an integration over the cumulative
distribution $g$ of absorption coefficients. That distribution is smooth and can be sampled at a few points
\citep{mlawer1997}. RRTMG LW uses 16 spectral bands (\code{nbndlw}) and 140 g-points in total (\code{ngptlw}). Each
band has its own absorbers:
\begin{itemize}
  \item water vapor, carbon dioxide and ozone;
  \item the minor gases N$_2$O, CH$_4$, CFCs;
  \item the water vapor continuum.
\end{itemize}
The absorption coefficients come from tables precomputed from line-by-line calculations. They are tabulated as a
function of pressure, temperature and the binary-species parameter of the gas pair that dominates the band.

\paragraph{The calculation for one column.} \src{phys/module_ra_rrtmg_lw.F}{rrtmg_lw} runs these steps:
\begin{description}
  \item[\code{inatm}] prepares the layer properties: pressures, temperatures, the column amounts of each gas and the
    cloud properties. Above the model top it adds extra layers, from a standard profile, up to 0.1~hPa. For the case
    the column has $n_L = 109$ layers, of which 50 are added.
  \item[\code{cldprmc}] computes the cloud optical depths per g-point from the cloud water and ice paths and the
    particle sizes.
  \item[\code{setcoef}] computes, for every layer, the table indices and interpolation weights in pressure and
    temperature, and the Planck function terms.
  \item[\code{taumol}] calls \code{taugb1} \dots{} \code{taugb16}, one routine per band, which interpolate the
    absorption coefficients for every layer and g-point: the optical depths $\tau_{g,k}$.
  \item[\code{rtrnmc}] performs the radiative transfer: for each g-point, a downward then an upward sweep through the
    layers with the two-stream source function and the diffusivity approximation, summed over g-points to the fluxes
    and the heating rate.
\end{description}

The heating rate of layer $k$ is
\begin{equation}
  \left.\pd{T}{t}\right|_k = \frac{g}{c_p}\,\frac{F^{\mathrm{net}}_{k+1} - F^{\mathrm{net}}_{k}}{\Delta p_k},
\end{equation}
which the driver converts into a tendency of $\theta$.

\paragraph{Cloud overlap: McICA.} Partial cloudiness in a column could be treated by overlap assumptions that
multiply the work. RRTMG uses the Monte Carlo Independent Column Approximation \citep{pincus2003} instead:
\begin{itemize}
  \item each g-point sees one randomly generated cloudy or clear sub-column, consistent with the layer cloud fractions
    and the overlap rule;
  \item the random numbers come from the KISS generator (\code{kissvec}), seeded from the column's pressures, so the
    same column always gets the same sub-columns.
\end{itemize}
\code{mcica_subcol_lw} and \code{generate_stochastic_clouds} implement this.

\begin{algorithm}[ht]
\caption{One RRTMG longwave column, \code{RRTMG_LWRAD} (simplified).}
\label{alg:rrtmg}
\begin{algorithmic}[1]
\State column setup: layers from the model plus buffer layers above the top (\code{inatm}); gases from \code{read_CAMgases}
\State McICA sub-columns: cloud yes/no per g-point and layer (\code{mcica_subcol_lw}, \code{kissvec})
\State cloud optical depths per g-point (\code{cldprmc})
\State table indices and Planck terms per layer (\code{setcoef})
\For{band $b = 1, \dots, 16$}
  \State gas optical depths for the band's g-points and all layers (\code{taugb}$b$)
\EndFor
\For{g-point $g = 1, \dots, 140$} \Comment{in band order}
  \State downward sweep $k = n_L, \dots, 1$; upward sweep $k = 1, \dots, n_L$ (\code{rtrnmc})
  \State add the g-point's fluxes to the band and total fluxes
\EndFor
\State heating rates and surface fluxes; copy back the model's layers
\end{algorithmic}
\end{algorithm}

\paragraph{Precision.} In this build RRTMG LW computes in single precision: \code{kind_rb = kind(1.0)}. The
double-precision branches of the source are switched off by a preprocessor condition.

\begin{portnote}
RRTMG is the largest single item of the physics port.
\begin{itemize}
  \item \textbf{Tables.} Its absorption tables (about 0.75~MB) live in modules, partly as \code{EQUIVALENCE}d pairs of
    arrays. The device cannot share storage that way, so the port flattens each pair into one array with explicit
    index arithmetic. The storage order is the same and so are the values. This is a shared refactor, proven bitwise
    on the CPU first.
  \item \textbf{Initialization.} \code{rrtmg_lw_ini} rebuilt the same tables at every radiation call. It is moved to
    the model's initialization, which is also bit-neutral.
  \item \textbf{Memory.} A column needs about 2{,}500$\,n_L$ words of temporaries, about 1.1~MB at $n_L = 109$. The
    GPU processes batches of columns (4096 by default) with one thread per column, and the temporaries become slices
    of batch work arrays.
  \item \textbf{Sums.} The fluxes are sums over g-points. A faster version that parallelizes over g-points must still
    add the contributions in the original order, or it changes the bits (\cref{ch:colphys}).
\end{itemize}
\end{portnote}

\section{Dudhia shortwave}

The scheme of \citet{dudhia1989} integrates the downward solar flux from the top of the atmosphere to the ground in a
single broad band, for each column independently (\src{phys/module_ra_sw.F}{SWPARA}). In each layer, the flux is
reduced by:
\begin{itemize}
  \item clear-air scattering, proportional to the layer's pressure depth (\code{cssca}) and divided by the cosine of
    the zenith angle $\mu_0$ (\code{XMU});
  \item water-vapor absorption, from the absorption of the path $u$ above the layer:
    \begin{equation}
      A(u) = \frac{2.9\,u}{(1 + 141.5\,u)^{0.635} + 5.925\,u}, \qquad u \text{ the vapor path in g\,cm}^{-2} \text{ over } \mu_0;
    \end{equation}
  \item cloud reflection and absorption, from two small tables (\code{ALBTAB}, \code{ABSTAB}) indexed by the
    logarithm of the liquid water path and by $\mu_0$, interpolated bilinearly.
\end{itemize}
The heating of each layer is the absorbed part of the flux; the ground receives what is left. Night columns
($\mu_0 \le 0$) skip the integration.

\begin{portnote}
\begin{itemize}
  \item The Dudhia scheme is a classic \emph{single column with an early exit}: in a batch of columns, the night ones
    finish at once while the day ones run the full loop. On a GPU this leaves threads of a warp idle (divergence,
    \cref{ch:warp}).
  \item Its \code{DATA} tables become \code{PARAMETER} arrays, so the device has them as constants.
  \item The power $0.635$ and the logarithm go through \code{rp_pow} and \code{rp_log10}.
\end{itemize}
\end{portnote}

\section{Ozone and greenhouse gases}

With \code{o3input = 2}, ozone comes from a climatology of monthly zonal means on 59 pressure levels. The driver
interpolates it in time to the model date and in pressure to the model layers. The pressure interpolation in
\code{ozn_p_int} searches, for each model level, the bracketing climatology levels. The CPU code shares the search
across the points of a $j$ row; the GPU version searches per column and finds the same bracket, so the weights and
results are identical (test T-OZN). Greenhouse-gas concentrations (CO$_2$, N$_2$O, CH$_4$, CFC-11, CFC-12) come from a
table by year (\code{ghg_input = 1}, \code{read_CAMgases}).
